\section{De DQN a Actor-Critic}

\subsection{Limitaciones de Q-Learning}

\begin{warningbox}{Dificultades de Q-Learning en espacios de acciones continuos}
Q-Learning necesita calcular $\max_{a'} Q(s', a')$. En un espacio de acciones discreto se pueden recorrer todas las acciones. Pero en un espacio de acciones continuo, esta maximización se convierte en un problema de optimización continua, con un costo computacional muy alto.
\end{warningbox}

\subsection{Métodos Actor-Critic}

Solución: aprender al mismo tiempo la política (actor) y la función de valor (critic):
\begin{itemize}
    \item \textbf{Actor}: $\pi_\theta(a|s)$, política parametrizada que produce las acciones.
    \item \textbf{Critic}: $Q_\phi(s, a)$, evalúa qué tan buenas son las acciones que elige el actor.
    \item El actor se actualiza con gradiente de política y el critic, con aprendizaje TD.
\end{itemize}

\subsection{Resumen de la sección}

Los métodos Actor-Critic combinan las ventajas de los métodos de gradiente de política y de los métodos basados en funciones de valor, y son el paradigma dominante para tratar espacios de acciones continuos.

